% Part: sets-functions-relations
% Chapter: infinite
% Section: dedekinds-proof
%
\documentclass[../../../include/open-logic-section]{subfiles}

\begin{document}

\olfileid{sfr}{infinite}{dedekindsproof}

\olsection[Dedekinds ``bewijs'']{Dedekinds ``bewijs'' voor het
bestaan van een oneindige verzameling}

In dit hoofdstuk hebben we de natuurlijke getallen
verzamelingtheoretisch behandeld met behulp van Dedekindalgebra's. In
\olref[arith][ref]{sec} stonden we onder meer stil bij de filosofische
betekenis van de aritmetisering van de analyse. Nu moeten we stilstaan
bij de betekenis van wat we hier hebben bereikt.

In heel \olref[sfr][arith][]{chap} namen we de natuurlijke getallen als
gegeven aan en construeerden we daaruit expliciet de gehele, rationale
en reële getallen. In dit hoofdstuk hebben we geen expliciete
constructie van de natuurlijke getallen gegeven. We hebben slechts
aangetoond dat we, \emph{uitgaande van een willekeurige
Dedekind-oneindige verzameling}, een verzameling kunnen definiëren die zich precies zo
gedraagt als~$\Nat$ zich volgens ons moet gedragen.

Daarmee kunnen we uiteraard niet beweren dat we een metafysische vraag
hebben beantwoord, zoals \emph{welke objecten zijn de natuurlijke
getallen}. Maar dat is maar goed ook. In
\olref[sfr][arith][ref]{sec} benadrukten we immers dat het onjuist zou
zijn om de definitie van~$\Real$ als de verzameling van Dedekindsneden
als een \emph{ontdekking} te beschouwen en niet als een handige afspraak.
Doorslaggevend is dat Dedekindsneden de belangrijkste wiskundige
eigenschappen van de reële getallen vertonen. Hetzelfde geldt hier:
\emph{iedere} Dedekindalgebra vertoont de belangrijkste wiskundige
eigenschappen van de natuurlijke getallen. (Dedekind onderstreepte dit
door te bewijzen dat alle Dedekindalgebra's \emph{isomorf} zijn
(\citeyear[stellingen 132--3]{Dedekind1888}). Het is dan ook niet
verwonderlijk dat veel hedendaagse ``structuralisten'' Dedekind als
voorloper aanhalen.)

Bovendien hebben we laten zien hoe de theorie van de natuurlijke
getallen kan worden ingebed in een na\"ieve, eenvoudige
verzamelingenleer. Die is zelf nog tamelijk informeel, maar
veronderstelt de natuurlijke getallen (naar het schijnt) niet als
gegeven. Mogelijk zijn we dus op weg om Dedekinds eigen ambitieuze
project te verwezenlijken. Hij legde dat als volgt uit:
\begin{quote}
	In de wetenschap mag niets wat kan worden bewezen, zonder bewijs
	worden geloofd. Hoewel deze eis redelijk lijkt, kan ik hem niet als
	vervuld beschouwen, zelfs niet bij de nieuwste methoden om de
	grondslagen te leggen van de eenvoudigste wetenschap, namelijk dat
	deel van de logica dat de getallentheorie behandelt. Met mijn
	uitspraak dat rekenkunde (algebra, analyse) slechts een deel van de
	logica is, bedoel ik dat ik het getalbegrip geheel onafhankelijk acht
	van de begrippen of intuïties van ruimte en tijd---dat ik het veeleer
	beschouw als een onmiddellijk voortbrengsel van de zuivere wetten van
	het denken.
	\citep[voorwoord]{Dedekind1888}
\end{quote}
Dedekinds gedurfde idee is dus het volgende. We hebben zojuist laten
zien hoe de natuurlijke getallen uitsluitend met (na\"ieve)
verzamelingenleer kunnen worden opgebouwd. In
\olref[sfr][arith][]{chap} zagen we hoe we de reële getallen kunnen
construeren als de natuurlijke getallen en enige verzamelingenleer
gegeven zijn. Misschien blijkt ``rekenkunde (algebra, analyse)'' dus
``slechts een deel van de logica'' te zijn (in Dedekinds verruimde
betekenis van het woord ``logica'').

Tot zover het idee. Maar er is een probleem. Onze constructie van een
Dedekindalgebra (onze vervanger voor de natuurlijke getallen)
veronderstelt dat er een Dedekind-oneindige verzameling bestaat. (Zie
nogmaals
\olref[sfr][infinite][dedekind]{thm:DedekindInfiniteAlgebra}.) Als het
bestaan van een Dedekind-oneindige verzameling niet met ``logica'' of
``de zuivere wetten van het denken'' kan worden aangetoond, loopt het
project vast.

\emph{Kan} het bestaan van een Dedekind-oneindige verzameling dan met
``de zuivere wetten van het denken'' worden aangetoond? Dedekind deed de
volgende poging:
\begin{quote}
  Mijn eigen gedachtenwereld, dat wil zeggen de totaliteit~$S$ van alle
  dingen die voorwerp van mijn denken kunnen zijn, is oneindig. Als~$s$
  namelijk een element van~$S$ aanduidt, dan is de gedachte~$s'$ dat~$s$
  voorwerp van mijn denken kan zijn, zelf een element van~$S$. Als we
  deze gedachte beschouwen als een beeld~$\phi(s)$ van het element~$s$,
  dan is \dots~$S$ [Dedekind-]oneindig, zoals te bewijzen was.
	\citep[\S66]{Dedekind1888}
\end{quote}
Het is zeer verrassend dit aan te treffen midden in een boek dat
grotendeels uit uiterst nauwkeurige wiskundige bewijzen bestaat. Twee
opmerkingen zijn op hun plaats.

Ten eerste: dit ``bewijs'' is naar huidige maatstaven nauwelijks
wiskundig. Het gaat over psychologische objecten (gedachten), en
bovendien slechts over \emph{mogelijke} objecten.

Ten tweede: althans bij onze presentatie van Dedekindalgebra's heeft dit
``bewijs'' een duidelijk technisch gebrek. Als Dedekinds redenering
slaagt, toont zij alleen aan dat er oneindig veel dingen zijn (meer
bepaald: oneindig veel gedachten). Maar Dedekind moet ons ook een reden
geven om~$S$ op te vatten als één \emph{verzameling} met oneindig veel
elementen, in plaats van~$S$ op te vatten als \emph{een aantal
dingen} (in het meervoud).

Dat Dedekind hier geen leemte zag, kan erop wijzen dat zijn gebruik van
het woord ``totaliteit'' niet precies overeenkwam met \emph{ons} gebruik
van het woord ``verzameling''.\footnote{Er zijn inderdaad andere redenen
om te denken dat dit niet zo was; zie \citet[p.~23]{Potter2004}.} Dat
hoeft niet erg te verbazen. Het project van de afgelopen twee
hoofdstukken---een ``constructie'' van de natuurlijke getallen en
vandaaruit een ``constructie'' van de gehele, reële en rationale
getallen---hebben we geheel na\"ief uitgevoerd. Telkens wanneer we een
verzameling nodig hadden, namen we eenvoudig aan dat die bestond, zonder
veel algemene beginselen vast te leggen over welke verzamelingen precies
bestaan en wanneer. We weten echter dat we \emph{enkele} algemene
beginselen nodig hebben; anders lopen we tegen de paradox van Russell
aan.

Het is tijd onze na\"iviteit achter ons te laten.
\end{document}